\begin{vidu}
Một dao động điều hoà trên đoạn thẳng dài 10 cm và thực hiện được 50 dao động trong thời gian 78,5 s. Lấy $\pi=3,14$.
\begin{itemize}
\item [1.] Biên độ của dao động là\dotfill
\item [2.] Chu kì của dao động là\dotfill
\item [3.] Tần số góc của dao động là\dotfill
\item [4.] Vận tốc của vật tại li độ 2,5 cm là\dotfill
\item [5.] Li độ của vật khi vận tốc có giá trị $10\sqrt{2}\ cm/s$ là\dotfill
\item [6.] Vận tốc của vật khi đi qua vị trí có li độ $x=-3\ cm$ theo chiều hướng về vị trí cân bằng là\dotfill
\item [7.] Gia tốc của vật khi đi qua vị trí có li độ $x=-3\ cm$ theo chiều hướng về vị trí cân bằng là\dotfill
\end{itemize}
\loigiai{
\begin{itemize}
\item $A=5\ cm;\ T=1,57\ s\Rightarrow \omega=4\ rad/s$.
\item $|v|=\omega\sqrt{A^2-x^2}=10\sqrt{3}\ cm/s$.
\item $|x|=\sqrt{A^2+\dfrac{v^2}{\omega^2}}=\dfrac{5\sqrt{2}}{2}\ cm/s$. 
\item $v=\omega \sqrt{A^2-x^2}=16\ cm/s;\quad a=-\omega^2 x=48\ cm/s^2$.

\end{itemize}
}
%\haidong{8}
\end{vidu} %\dotlineEX{1}
%\loai{Cho khoảng cách từ vận tới vị trí cân bằng, tìm v}
\begin{vidu} %[3P1-2.2B][Loai1]  
Một vật dao động điều hoà với biên độ 6 cm, tần số góc 4$\pi$ rad/s. Tốc độ của vật tại điểm cách vị trí cân bằng một đoạn 2 cm là
\choice
{6$\sqrt{12}$ cm/s}
{22 cm/s}
{\True 16$\pi \sqrt{2}$ cm/s}
{16$\pi$ cm/s}
\loigiai{
\begin{itemize}
\item Theo phương trình độc lập thời gian: $\dfrac{x^2}{A^2} + \dfrac{v^2}{\omega^2.A^2} = 1$.
\item Thay số ta được: $\dfrac{2^2}{6^2} + \dfrac{v^2}{\left(4\pi\right)^2.6^2} = 1\Rightarrow |v| = 16\pi.\sqrt{2}$ cm/s.
\end{itemize}
}
%\haidong{3}
\end{vidu} %\dotlineEX{2}
%\loai{Cho li độ x tìm v}
\begin{vidu} %[3P1-2.2B][Loai1]  
Cho một chất điểm dao động điều hòa với tần số 3 Hz. Tốc độ cực đại trong quá trình dao động bằng 30$\pi$ cm/s. Khi chất điểm đi ngang qua vị trí có li độ 4 cm, tốc độ chuyển động của chất điểm bằng
\choice
{8$\pi$ cm/s}
{\True 18$\pi$ cm/s}
{6$\pi \sqrt{3}$ cm/s}
{4$\pi \sqrt{3}$ cm/s}
\loigiai{
\begin{itemize}
\item Ta có: $\omega = 2\pi.3 = 6\pi \Rightarrow A = \dfrac{v_{max}}{\omega} = \dfrac{30\pi}{6\pi} = 5\ cm$.
\item Áp dụng phương trình độc lập thời gian, ta có: $\dfrac{4^2}{5^2} + \dfrac{v^2}{\left(30\pi\right)^2} = 1\Rightarrow |v| = 18\pi$ cm/s.
\end{itemize}
}
%\haidong{4}
\end{vidu} %\dotlineEX{2}
%\loai{Dùng công thức độc lập x,v tìm $\omega$, T, f}
\begin{vidu}%[3P1-2.2B][Loai1]
Một vật dao động điều hòa trên trục Ox. Khi vật qua vị trí cân bằng, tốc độ của nó là $8\pi$  cm/s. Khi vật cách vị trí cân bằng 3,2 cm thì nó có tốc độ là $4,8\pi$  cm/s. Tần số của dao động là
\choice
{4 Hz}
{0,5 Hz}
{2 Hz}
{\True 1 Hz}
\loigiai{
\begin{itemize}
\item Ta có: $v_{\max} = 8\pi$  cm/s.
\item Luôn có: ${{\left( \dfrac{x}{A} \right)}^2}+{{\left( \dfrac{v}{v_{1\max}} \right)}^2}=1\Rightarrow \left(\dfrac{3,2}{A} \right)^2+\dfrac{\left(4,8\pi\right)^2}{\left(8\pi \right)^2}=1^2\Rightarrow A=4\ cm \Rightarrow \omega  = 2\pi\ rad/s  \Rightarrow f = 1\ Hz$.
\end{itemize}
}
%\haidong{4}
\end{vidu} %\dotlineEX{2}

%\loai{Dùng công thức độc lập x,v tìm $v_{\max}$}
\begin{vidu}%[3P1-2.2K][Loai1]
Một dao động điều hòa có vận tốc và tọa độ tại thời điểm $t_1$ và $t_2$ tương ứng là: $v_1 = 20$ cm/s; $x_1 = 8\sqrt{3}$ cm và $v_{2} = 20\sqrt{2}$ cm/s ; $x_2 = 8\sqrt{2}$ cm. Vận tốc cực đại của dao động là
\choice
{$40\sqrt{2}$ cm/s}
{80 cm/s}
{\True 40 cm/s}
{$40\sqrt{3}$ cm/s}
\loigiai{
\begin{itemize}
\item Ta có: 
$\begin{cases}
x_1^2 + \dfrac{v_1^2}{\omega^2} = A^2\quad (1)\\
x_2^2 + \dfrac{v_2^2}{\omega^2} = A^2\quad (2)
\end{cases} 
\Rightarrow x_1^2 + \dfrac{v_1^2}{\omega^2} = x_2^2 + \dfrac{v_2^2}{\omega^2} \Rightarrow \omega = \sqrt {\dfrac{v_1^2 - v_2^2}{x_2^2 - x_1^2}} = 2,5\ rad/s$. 
\item Biên độ dao động: $ A^2=x_1^2+\dfrac{v_1^2}{\omega^2} \Rightarrow A=16\ cm\Rightarrow v_{\max} = A\omega  = 40$ cm/s.

\end{itemize}
}
%\haidong{5}
\end{vidu} %\dotlineEX{2}
%\loai{Công thức độc lập giữa a và v}
\begin{vidu}%[3P1B2-2][Loai1]
Vật dao động điều hòa trên trục Ox. Khi vật qua vị trí cân bằng có tốc độ 20 cm/s. Khi vật có tốc độ 10 cm/s thì độ lớn gia tốc của vật là $50\sqrt{3}\ cm/s^2$. Biên độ dao động của vật là
\choice
{5 cm}
{\True 4 cm}
{3 cm}
{2 cm}
\loigiai{
\begin{itemize}
\item Khi qua vị trí cân bằng, vật có tốc độ cực đại $v_{1\max}=\omega A=20\ cm/s$.
\item $\dfrac{v^2}{\left(\omega A\right)^2}+\dfrac{a^2}{\left(\omega^2A\right)^2}=1\Leftrightarrow \dfrac{10^2}{20^2}+\dfrac{\left( 50\sqrt{3} \right)^2}{\left( 20\omega \right)^2}=1 \Rightarrow \omega =5\ rad/s \Rightarrow A = 4$ cm.
\end{itemize}
}
%\haidong{5}
\end{vidu} %\dotlineEX{2}
%\loai{Công thức độc lập giữa a và x}
\begin{vidu}%[3P1-2.2K][Loai1]
Một vật dao động điều hòa, tại vị trí có li độ $- 1$ cm thì gia tốc là 1 $m/s^2$. Tại vị trí có li độ 4 cm độ lớn gia tốc bằng
\choice
{$- 4\ m/s^2$}
{\True 4 $m/s^2$}
{8 $m/s^2$}
{2 $m/s^2$}
\loigiai{
\begin{itemize}
\item Ta có: $a = - \omega^2x$, khi $x = - 1$ cm thì $a = 1\ m/s^2 = 100\ cm/s^2 \Rightarrow \omega^2 = 100$.
\item Khi $x = 4$ cm thì $a = - 100.4 = -400\ cm/s^2 = - 4\ m/s^2$.
\item Độ lớn gia tốc là 4 $m/s^2$.
\end{itemize}
}
%\haidong{4}
\end{vidu} %\dotlineEX{3}
%\loai{Cho hai gia tốc. Tìm gia tốc tại vị trí thứ 3}
\begin{vidu}%[3P1K2-2][Loai1]
Một chất điểm dao động điều hoà trên một đoạn thẳng, khi đi qua M và N trên đoạn thẳng đó chất điểm có gia tốc lần lượt là $a_M = - 3\ m/s^2$ và $a_N = 6\ m/s^2$. C là một điểm trên đoạn MN và $CM = 2CN$. Gia tốc chất điểm khi đi qua C có giá trị là
\choice
{1 $m/s^2$}
{2 $m/s^2$}
{\True 3 $m/s^2$}
{4 $m/s^2$}
\loigiai{
\begin{itemize}
\item Ta có: $a=-\omega^2x$, nghĩa rằng a và x luôn trái dấu nhau $\Rightarrow$ M có li độ dương và N có li độ âm ($x_M > x_N$).
\item Lại có: $CM = 2CN$, do đó: $x_M-x_C=2x_C-x_N\Leftrightarrow 3x_C=x_M+2x_N\Rightarrow x_C=\dfrac{x_M+2x_N}{3}$.
\item Mà $a=-\omega^2x\Rightarrow a\sim x$ (a tỉ lệ với x) $\Rightarrow$ $a_C=\dfrac{a_M+2a_N}{3}=\dfrac{-3+2.6}{3}=3\ m/s^2$. 

\end{itemize}
}
%\haidong{6}
\end{vidu} %\dotlineEX{6}
\begin{comment}%[3P1KD-1][Loai1]
Một vật dao động điều hòa với biên độ A trên trục Ox. Xét quá trình vật đi từ vị trí cân bằng ra biên, khi vật rời khỏi vị trí cân bằng đoạn S thì tốc độ của vật là $A\sqrt{8}$ m/s, vật đi thêm đoạn S nữa thì tốc độ giảm còn $A\sqrt{5}$ m/s, vật đi thêm đoạn S nữa thì tốc độ của vật là bao nhiêu? Biết $3S \leq A$.
\choice
{2A m/s}
{\True 0}
{3A m/s}
{A m/s}
\loigiai{
\begin{itemize}
\item Sử dụng công thức độc lập thời gian x và v cho các vị trí ta có: 
\begin{align*}
\begin{cases}
S^2 + \dfrac{8A^2}{\omega^2} = A^2\\ 
4S^2 + \dfrac{5A^2}{\omega^2} = A^2 
\end{cases} 
\Rightarrow \begin{cases}A = 3S\\ \omega =3\end{cases} \end{align*}
\item Khi vật đi thêm đoạn S nữa: $9S^2 + \dfrac{v^2}{\omega^2} = A^2  \Rightarrow v = 0\quad  \ct{(vật tới biên sau khi đi 3S)}$.

\end{itemize}
}
%\haidong{12}
\end{comment} %\dotlineEX{9}
